\subsection{One query budget for several comparisons}
An auditor comparing several models can reuse each reading across their
loss contrasts. Let $W_{r,Q}(y)$ be the sharp width for contrast $r$ after
the common transcript $Q$, and let $W_{r,0}>0$ be its initial width. Define
\[
 J_Q(y)=\max_r W_{r,Q}(y)/W_{r,0}.
\]
Then $J_Q(y)\leq1/2$ means that every requested comparison has reached
half its original width. Already determined contrasts can be omitted.
This is the local diameter of the vector contrast map in a weighted
$\ell_\infty$ norm, within the established vector-recovery framework
\citep{krieg2025}; shared-label model certification also has close prior
work \citep{kuboyama2026}.

The scalar assumption in Theorem~\ref{prop:risk_classification} matters:
the maximum of concave scalar widths need not be concave. In the same
six-interval prefix interface, two constructed contrasts have optimal
randomized-adaptive mean joint width $1/2$, versus $2/3$ for the optimal
randomized batch with two queries. Online Resource~1, S8 gives the
coefficients, policies, matching adverse witnesses and complete proof.
These are mean-width values, not precision probabilities or measured savings.

For the uncapped exact-total simplex, each batch's scalar widths are
linear in the energy allocation. Thus $J_Q$ is convex. Its worst mixture
expectation is attained at a concentration vertex, making the finite-game
LP exact for randomized-batch mean risk. The deterministic two-query
planner instead minimizes the worst joint width after feedback. Its
conditional objective is the maximum of the contrast-specific affine
widths; checking all line intersections gives the exact optimum.
The two objectives remain distinct. A deterministic adaptive value below
the randomized-batch lower bound establishes an adaptive advantage; the
reverse inequality does not establish optimality among adaptive policies.
The vertex argument does not transfer to precision indicators, caps or
measurement error. S8 provides the formulations and verification details.
